\documentclass[final,a4]{beamer}
\usetheme{Madrid}
\title{le projet RefPerSys}
\author{Basile STARYNKEVITCH}
\date{mai 2025}
\usepackage{relsize}
\usepackage{luacode}
\usepackage{xcolor}
\usepackage{alltt}
%\usepackage{wasysym}
\usepackage{hyperref}
%\usepackage{newunicodechar}

%- \hypersetup{
%-   colorlinks   = true, %Colours links instead of ugly boxes
%-   urlcolor     = NavyBlue, %Colour for external hyperlinks
%-   linkcolor    = DarkGreen, %Colour of internal links
%-   citecolor   = DarkMagenta, %Colour of citations
%-   frenchlinks = true,
%- }

\begin{document}
 \begin{luacode*}
   local gitpip=io.popen("git log --no-color --format=oneline -1 --abbrev=16 --abbrev-commit -q | cut -d' ' -f1")
   gitid=gitpip:read()
   gitpip:close()
 \end{luacode*}
 \newcommand{\mygitid}{\luadirect{tex.print(gitid)}}

 \begin{frame}
   \titlepage

   \small{(un moteur d'inférences, ``intelligence artificielle'' symbolique)}

   \bigskip

   \begin{center}
   \begin{relsize}{-2}
     8 rue de la Faïencerie\\
     92340 Bourg-la-Reine
   \end{relsize}

   \bigskip

   \textbf{les opinions me sont personnelles}
   
   \bigskip
   
     \texttt{basile{@}starynkevitch.net}

     \bigskip

     \begin{fbox}
       {Journées Du Logiciel Libre - Lyon - mai 2025}
     \end{fbox}
     

     git {\texttt \mygitid}
   \end{center}
   
 \end{frame}

 \begin{frame}{Plan}

   \begin{relsize}{-1.5}
   \tableofcontents
   \end{relsize}
   
   \bigskip

   \begin{center}
     \relsize{+2}{\href{http://refpersys.org/}{\texttt{refpersys.org}}}

     \bigskip
     \includegraphics[width=0.45\textwidth]{refpersys-logo}

     \smallskip

     \relsize{-2}{Logo par Gaëtan Tapon}
   \end{center}
 \end{frame}

 \section{Introduction}
 \label{sec:intro}
 
 \begin{frame}{motivations pour RefPerSys}

   \begin{itemize}
   \item souveraineté numérique \small{(moins dépendre de logiciels
     états-uniens)}
   \item s'occuper 
   \item expérimenter un logiciel
   \item populariser des idées méconnues (celles de Jacques Pitrat)
   \item \textbf{un projet \textcolor{red}{en cours de dévelopement} qui cherche des contributeurs}
   \end{itemize}

   \bigskip
   \textbf{Le code source est libre} sous licence GPLv3+/CeCILLv2
   \textbf{sans garantie} sur
   \href{https://github.com/RefPerSys/RefPerSys}{\texttt{github.com/RefPerSys/RefPerSys}}
   et peut-être même difficilement utilisable en mai 2025.

   \bigskip

   \begin{center}
     \large\textsc{\textcolor{red}{REF}lexive \textcolor{red}{PER}sistent \textcolor{red}{SYS}tem}
   \end{center}
 \end{frame}
 
 \begin{frame}{historique et systèmes inspirants pour RefPerSys}

   \begin{itemize}

   \item \textbf{les moteurs d'inférences des systèmes experts} des années 1990
     \href{https://clipsrules.net/}{\textsc{CLIPSrules}} (NASA)

   \item \textbf{le système \textcolor{red}{CAIA}} (``chercheur artificiel en intelligence artificielle'') de J.Pitrat
     
   \item \textbf{les outils d'aide à la démonstration} ou assistants
     de preuve \href{https://rocq-prover.org/}{\textsc{Rocq}} (ex
     \textsc{Coq})

     \item La problématique de \textbf{\textcolor{red}{l'amorçage des compilateurs}}: tous les
       compilateurs se compilent eux-même.
   \end{itemize}

   \bigskip

   En France, il existe une grosse culture mathématique de systèmes
   formels de preuve. Ma philosophie est plus pragmatique.
   
 \end{frame}
\end{document}

%%%%%%%%%%%%%%%%%%%%%%%%%%%%%%%%%%%%%%%%%%%%%%%%%%%%%%%%%%%%%%%%
%% Local Variables: ;;
%% compile-command: "./build-refpersys-jdll2025.sh" ;;
%% End: ;;
%%%%%%%%%%%%%%%%%%%%%%%%%%%%%%%%%%%%%%%%%%%%%%%%%%%%%%%%%%%%%%%%
